\documentclass[11pt]{article}
\usepackage[utf8]{inputenc}
\usepackage{amsmath,amssymb}
\usepackage[margin=1in]{geometry}
\usepackage{hyperref}
\emergencystretch=3em

\title{The fluctuation closed form: $S_\lambda(a)=2^{1-\lambda}\beta^{-\lambda}\Gamma(2\lambda)
e^{\beta a^2/8}D_{-2\lambda}(-a\sqrt{\beta/2})$}
\author{Claude, with Daniel Murfet}
\date{2026-09-06}

\begin{document}
\maketitle
\sloppy

\begin{abstract}
Unit 7 of the grammar-paper \S4 autoformalisation, and the payoff of the option-B parabolic-cylinder
detour: the corollary \texttt{cor:fluctuation\_closed\_form}, which had been the arc's standing blocker.
The change of variables $t=u^2/(2\beta)$ carries Watanabe's fluctuation function onto the
$D_{-2\lambda}$ integral kernel; the $\Gamma$ and Gaussian prefactors then cancel. Proved via an
intermediate $\Gamma$-free ``kernel'' lemma, so that all substitution bookkeeping stays disentangled
from the special-function definition. Zero \texttt{sorry}/\texttt{axiom}.
\end{abstract}

\section{Setting}
Unit~6 built the negative-order parabolic cylinder function \texttt{parCylNeg}
$=D_{-\nu}(x)=e^{-x^2/4}/\Gamma(\nu)\int_0^\infty u^{\nu-1}e^{-u^2/2-xu}\,du$ and its integrability.
This unit is the reason it was built: the paper's closed form for $S_\lambda$. With it, the entire
analytic core of \S4 that does not require an operator/commutator framework is formalised --- the
ladder, the recurrence, the ODE, and now the parabolic-cylinder identification.

\section{The things formalised}
{\small
\begin{verbatim}
theorem fluctuation_eq_kernel (β lam a : ℝ) (hβ : 0 < β) :
    fluctuation β lam a
      = (2 ^ (1 - lam) * β ^ (-lam))
        * ∫ u in Set.Ioi 0,
            u ^ (2*lam - 1) * Real.exp (-u^2/2 + a * Real.sqrt (β/2) * u)

theorem fluctuation_closed_form (β lam a : ℝ) (hβ : 0 < β) (hlam : 0 < lam) :
    fluctuation β lam a
      = 2 ^ (1 - lam) * β ^ (-lam) * Real.Gamma (2*lam) * Real.exp (β*a^2/8)
        * parCylNeg (2*lam) (-a * Real.sqrt (β/2))
\end{verbatim}
}

\section{Proof strategy}
The substitution $t=u^2/(2\beta)$ is done in two library steps rather than one square-root change of
variables, following GPT-6 Astra's steer. First scale $t\mapsto r/(2\beta)$ with
\texttt{integral\_comp\_mul\_left\_Ioi} (Jacobian $(2\beta)^{-1}$); then apply
\texttt{integral\_comp\_rpow\_Ioi\_of\_pos} \emph{backwards} at $p=2$ (Jacobian $2u$). Their product,
$(2\beta)^{-1}\cdot 2u=u/\beta$, is exactly $dt/du$. A single pointwise identity on $(0,\infty)$ then
matches the transported integrand
$2u\cdot F(u^2/(2\beta))=2(2\beta)^{-(\lambda-1)}\,u^{2\lambda-1}e^{-u^2/2+a\sqrt{\beta/2}\,u}$,
and \texttt{integral\_const\_mul} pulls the constant out. Collecting the two Jacobians and the
$(2\beta)^{-(\lambda-1)}$ against the target prefactor is the scalar identity
$(2\beta)^{-1}\!\cdot 2(2\beta)^{-(\lambda-1)}=2^{1-\lambda}\beta^{-\lambda}$
(\texttt{hconstant}). This is \texttt{fluctuation\_eq\_kernel}, which never mentions $\Gamma$.

The corollary then unfolds \texttt{parCylNeg}$(2\lambda,-a\sqrt{\beta/2})$: its kernel integrand
$e^{-u^2/2-(-a\sqrt{\beta/2})u}$ is (after one \texttt{ring} rewrite) the kernel of
\texttt{fluctuation\_eq\_kernel}, its Gaussian prefactor is $e^{-\beta a^2/8}$ (since
$(-a\sqrt{\beta/2})^2=\beta a^2/2$), and its $1/\Gamma(2\lambda)$ cancels the $\Gamma(2\lambda)$ in the
target. \texttt{field\_simp} closes it with $\Gamma(2\lambda)\ne0$ and $e^{\beta a^2/8}\ne0$.

\section{Roadblocks and resolutions}
The route was Astra-planned, so the friction was purely Lean mechanics:
\begin{itemize}
\item The two change-of-variables lemmas live in \texttt{namespace MeasureTheory} inside a plain
\texttt{section IoiChangeVariables} (grep for \texttt{namespace} finds only the enclosing one); with
\texttt{open MeasureTheory} the unqualified names resolve.
\item The pointwise power identity splits into three separately-stated \texttt{have}s
(\texttt{hpow}, \texttt{hexp}, \texttt{hxx}): \texttt{ring} cannot see through \texttt{rpow}, so each
exponent identity ($2(\lambda-1)=2\lambda-2$, $1+(2\lambda-2)=2\lambda-1$) is normalised by an
explicit \texttt{show \ldots\ by ring} before \texttt{Real.rpow\_mul}/\texttt{rpow\_add}.
\item Two square-root facts are proved by \emph{avoiding} $(\sqrt{\cdot})^2$ under \texttt{field\_simp}:
$\sqrt{2\beta}\sqrt{\beta/2}=\beta$ via \texttt{Real.sqrt\_mul}$+$\texttt{sqrt\_sq}, and then
$\beta/\sqrt{2\beta}=\sqrt{\beta/2}$; substituting the latter backwards keeps the exponent-phase
identity rational in $\sqrt{2\beta}$ (never squared), so \texttt{field\_simp; ring} closes it.
\item \texttt{setIntegral\_congr\_fun} hands its pointwise goal as applied lambdas ---
\texttt{beta\_reduce} before the \texttt{rw}.
\item Several \texttt{field\_simp} calls closed their goals outright; the trailing \texttt{ring}s then
errored ``no goals'' and were deleted (a recurring repo pattern).
\end{itemize}

\section{What was Mathlib, what was new}
The change-of-variables engine is entirely Mathlib (\texttt{integral\_comp\_mul\_left\_Ioi},
\texttt{integral\_comp\_rpow\_Ioi\_of\_pos}, \texttt{integral\_const\_mul}, the \texttt{rpow}/\texttt{sqrt}
API). What is new is the pairing of the two Jacobians to realise the square-root substitution without a
square-root change-of-variables lemma, the $\Gamma$-free intermediate factorisation, and the
prefactor cancellation that identifies $S_\lambda$ with a named special function Mathlib does not have.

\section{Lessons}
A square-root change of variables is often cheaper as \emph{scale-then-square} than as one direct
substitution: the two off-the-shelf Jacobians multiply to the one you want, and each half is a lemma
Mathlib already proves unconditionally. And factoring the special-function definition out of the
substitution (prove the bare kernel integral first, unfold the definition only at the end) keeps
$\Gamma$ out of the change-of-variables algebra entirely --- the cancellation becomes a two-line
\texttt{field\_simp}.

\section{Follow-ups}
With the closed form landed, the D$_\nu$ arc continues to \texttt{eq:Dn\_hermite} (the integer-order
$D_n$ against Mathlib's Hermite polynomials) and \texttt{lem:qho} (the quantum harmonic oscillator
identification). The remaining non-special-function \S4 items --- \texttt{prop:fluctuation\_weber}
(ODE$\to$Weber transform) and the commutator $[b,b^\dagger]=1$ --- are independent of $D_\nu$.
\end{document}
